% Part: proof-theory
% Chapter: sequent-calculus
% Section: quantifiers

\documentclass[../../../include/open-logic-section]{subfiles}

\begin{document}

\olfileid[pt]{pt}{seq}{qua}

\olsection{Provas regulares e substituição}

As regras para os quantificadores introduzem algumas complicações por
causa das ``condições de variável própria'' que se aplicam às regras
\LeftR{\lexists} e~\RightR{\lforall}. Podemos tratar a maioria dessas
complicações garantindo que a !!{constant}~$a$ usada nessas inferências
nunca seja reutilizada. Introduzimos a seguinte definição.

\begin{defn}
!!^a{proof}~$\pi$ em \Log{G1c} ou \Log{G3c} é \emph{regular} se
\begin{enumerate}
  \item nenhuma variável própria~$a$ é usada como variável própria de
  mais de uma inferência \LeftR{\lexists} ou~\RightR{\lforall}; e
  \item se $a$ é a variável própria de uma inferência \LeftR{\lexists}
  ou~\RightR{\lforall}, ela ocorre apenas acima dessa inferência
  em~$\pi$.
\end{enumerate}
\end{defn}

\begin{lem}\ollabel{lem:subst-regular} Se $\pi$ é regular, termina em
$\Gamma \Sequent \Delta$, $c$ é !!a{constant} que não é usada como
variável própria em~$\pi$, e $t$~é um termo que não contém nenhuma
variável própria de~$\pi$, então $\Subst{\pi}{t}{c}$, obtida de~$\pi$
pela substituição de toda ocorrência de~$c$ por~$t$, é uma !!{proof}
regular. O sequente final de $\Subst{\pi}{t}{c}$ é
$\Subst{\Gamma}{t}{c} \Sequent \Subst{\Delta}{t}{c}$, onde
$\Subst{\Gamma}{t}{c} = \Setabs{!A(t)}{!A(c) \in \Gamma}$.
\end{lem}

\begin{proof}
  Procedemos por indução em $\pheight{\pi}$. Se $\pheight{\pi} = 0$,
  a prova consiste simplesmente em um axioma. Então
  $\Subst{\pi}{t}{c}$ também é um axioma, pois toda !!{formula}~$!A(c)$
  nela se transforma em $!A(t)$.

  Se $\pheight{\pi} > 0$, distinguimos casos conforme a última
  inferência em~$\pi$. Os casos das regras estruturais (em \Log{G1c})
  e das regras lógicas para os conectivos proposicionais são simples e
  ficam como exercícios. Consideramos os casos em que $\pi$ termina em
  uma inferência de quantificador.
  \begin{enumerate}
    \item Se a última inferência é $\RightR{\lforall}$, $\pi$ tem a
    forma
    \[
    \AxiomC{}
    \RightLabel{$\pi'$}
    \Deduce$\Gamma(c) \fCenter \Delta'(c), !A(a,c)$
    \RightLabel{\RightR{\lforall}}
    \UnaryInf$\Gamma(c) \fCenter \Delta'(c), \lforall[x][!A(x,c)]$
    \DisplayProof
    \]
    Pela hipótese de indução, $\Subst{\pi'}{t}{c}$ é uma !!{proof}
    regular de $\Gamma(t) \fCenter \Delta'(t), !A(a,t)$. Como $a$ é
    uma variável própria de~$\pi$, $a$~não ocorre em~$t$. Portanto, a
    última inferência em $\Subst{\pi}{t}{c}$,
    \[
    \AxiomC{}
    \RightLabel{$\Subst{\pi'}{t}{c}$}
    \Deduce$\Gamma(t) \fCenter \Delta'(t), !A(a,t)$
    \RightLabel{\RightR{\lforall}}
    \UnaryInf$\Gamma(t) \fCenter \Delta'(t), \lforall[x][!A(x,t)]$
    \DisplayProof
    \]
    é correta: a conclusão não contém~$a$. 
    \item Se a última inferência é $\LeftR{\lforall}$ em \Log{G3c},
    $\pi$ tem a forma
    \[
    \AxiomC{}
    \RightLabel{$\pi'$}
    \Deduce$!A(s(c),c), \lforall[x][!A(x,c)], \Gamma(c) \fCenter \Delta'(c)$
    \RightLabel{\RightR{\lforall}}
    \UnaryInf$\lforall[x][!A(x,c)], \Gamma(c) \fCenter \Delta'(c)$
    \DisplayProof
    \]
    Pela hipótese de indução, $\Subst{\pi'}{t}{c}$ é uma !!{proof}
    regular de $!A(s(t),t), \lforall[x][!A(x,t)], \Gamma(t) \fCenter
    \Delta'(t)$. A última inferência em $\Subst{\pi}{t}{c}$,
    \[
    \AxiomC{}
    \RightLabel{$\Subst{\pi'}{t}{c}$}
    \Deduce$!A(s(t),t), \lforall[x][!A(x,t)], \Gamma(t) \fCenter \Delta'(t)$
    \RightLabel{\RightR{\lforall}}
    \UnaryInf$\lforall[x][!A(x,t)], \Gamma(t) \fCenter \Delta'(t)$
    \DisplayProof
    \]
    é correta. O caso de \LeftR{\lforall} em~\Log{G1c} é semelhante,
    porém um pouco menos complicado.
    \item A última inferência é \RightR{\lexists} ou \LeftR{\lexists}:
    exercício.
  \end{enumerate}
  Em cada caso, acrescentamos um sequente ao final de uma !!{proof}
  regular (ou de duas !!{proof}s regulares, no caso de regras com duas
  premissas). O novo sequente difere do sequente final de $\pi$ apenas
  porque $c$ foi substituída pelo termo~$t$. Além disso, as variáveis
  próprias de $\pi$ e, como supusemos que $t$ não contém nenhuma
  variável própria de~$\pi$, as de $\Subst{\pi}{t}{c}$ são as mesmas; o novo
  sequente final não contém nenhuma variável própria de
  $\Subst{\pi}{t}{c}$. Consequentemente, $\Subst{\pi}{t}{c}$ é regular.
\end{proof}

\begin{prop}
Se $\pi$ é !!a{proof} em \Log{G1c} ou \Log{G3c}, então existe uma
!!{proof} regular~$\pi'$ com a mesma estrutura (em particular, a mesma
!!{height}) que difere de $\pi$ apenas pela substituição de variáveis
próprias.
\end{prop}

\begin{proof}
Somente para os fins desta prova, chamemos uma inferência com variável
própria em~$\pi$ de \emph{limpa} se sua variável própria não é a
variável própria de nenhuma outra inferência nem ocorre em qualquer
outro lugar de~$\pi$, exceto na sub-!!{proof} imediata da inferência;
caso contrário, chamemo-la de \emph{suja}. Seja $n$ o número de
inferências sujas com variável própria em~$\pi$. Mostramos, por indução
em~$n$, que $\pi$~pode ser transformada na !!{proof} regular exigida
~$\pi'$. Se $n=0$, todas as inferências com variável própria em~$\pi$
são limpas, de modo que $\pi$ é regular e nada há a provar. Caso
contrário, escolha uma inferência superior com variável própria, digamos
\RightR{\lforall}. A sub-!!{proof}~$\pi_1$ de~$\pi$ que termina nela tem
a forma
    \[
    \AxiomC{}
    \RightLabel{$\pi_2$}
    \Deduce$\Gamma \fCenter \Delta, !A(c)$
    \RightLabel{\RightR{\lforall}}
    \UnaryInf$\Gamma \fCenter \Delta, \lforall[x][!A(x)]$
    \DisplayProof
    \]
Observe que, como $c$ é uma variável própria, ela não ocorre
em~$\Gamma$, $\Delta$ ou $\lforall[x][!A(x)]$. A prova $\pi_2$ é
regular, pois não contém nenhuma inferência com variável própria. Seja
$a$ !!a{constant} que não ocorre em~$\pi$. Por
\olref{lem:subst-regular}, $\Subst{\pi_2}{a}{c}$ é uma prova regular de
$\Gamma \fCenter \Delta, !A(a)$, à qual podemos aplicar
~$\RightR{\lforall}$. Se substituirmos $\pi_1$ em~$\pi$ pela prova
$\pi_1'$,
    \[
    \AxiomC{}
    \RightLabel{$\Subst{\pi_2}{a}{c}$}
    \Deduce$\Gamma \fCenter \Delta, !A(a)$
    \RightLabel{\RightR{\lforall}}
    \UnaryInf$\Gamma \fCenter \Delta, \lforall[x][!A(x)]$
    \DisplayProof
    \]
substituímos uma inferência suja com variável própria por uma limpa, e a
única diferença é a substituição da variável própria $c$ por~$a$. A
prova resultante tem $n-1$ inferências sujas com variável própria, e o
resultado decorre da hipótese de indução.
\end{proof}

Não recorreremos explicitamente à proposição que acabamos de provar,
mas suporemos que todas as !!{proof}s consideradas são regulares---se
não forem, sempre podemos substituir as variáveis próprias para torná-las
regulares.

\end{document}
